\documentclass{article}

\usepackage{neurips_2026}
\usepackage[utf8]{inputenc}
\usepackage[T1]{fontenc}
\usepackage{hyperref}
\usepackage{url}
\usepackage{booktabs}
\usepackage{graphicx}
\usepackage{amsmath}
\usepackage{xcolor}
\usepackage{microtype}

\graphicspath{{figures/}}

\title{Anatomy-Guided Masking Does Not Help:\\
A Segmentation-Free Heuristic Outperforms a Trained Segmenter\\
for Joint-Embedding Pretraining on Retinal OCT}

\author{}

\begin{document}

\maketitle

\begin{abstract}
Masking policy is an under-audited design choice in medical self-supervised
learning. Intuitively, a masked-prediction model should be asked to reconstruct
diagnostically relevant tissue, so masking guided by a trained segmentation
model ought to beat unguided masking. We test that intuition directly. Starting
from a single shared checkpoint, we continue patch-level I-JEPA pretraining on
3D retinal OCT under six masking policies that differ only in how predictor
targets are placed, and evaluate all of them with one frozen linear probe
protocol on glaucoma classification (FairVision, $N{=}3000$).
The intuition fails. The best policy is an \emph{oracle} band located by a
first-order intensity statistic that consults no segmentation model
(AUC $0.8855$), beating the unguided null by $+0.0109$
(paired bootstrap 95\% CI $[+0.0058, +0.0162]$, $p<0.0005$) and beating a
policy driven by a trained retinal segmenter by $+0.0047$. Policies that place
targets most precisely on tissue --- $98\%$ of masked cells on anatomy ---
perform \emph{worst}. We measure mask geometry for every arm and show that
anatomical precision is not the operative variable, and we report a subgroup
analysis in which Black participants have the lowest AUC in every valid probe.
Encoder and probe weights are public and the headline number reproduces to
$10^{-5}$ across hardware and numerical precision.
\end{abstract}

\section{Introduction}

Self-supervised pretraining is increasingly applied to medical imaging, where
labels are scarce and anatomy is highly structured. I-JEPA~\citep{assran2023ijepa}
learns by predicting the representation of masked target blocks from a visible
context block. Which patches become targets is a free design choice, and in the
medical setting there is an obvious hypothesis: target the tissue that carries
the diagnosis.

We test that hypothesis on glaucoma classification from 3D optical coherence
tomography (OCT). All arms fork from one shared pretrained checkpoint at epoch
25 and differ \emph{only} in masking policy thereafter, so any downstream
difference is attributable to masking rather than to data, schedule, or
initialisation.

Our contributions:

\begin{itemize}
\item A controlled six-arm comparison of masking policies for medical SSL,
      evaluated under one frozen-probe protocol (21 probes total).
\item The finding that a segmentation-free intensity heuristic \emph{beats} a
      trained-segmenter-guided policy, and that increasing anatomical precision
      of the mask \emph{degrades} downstream AUC.
\item Direct measurement of mask geometry for every arm, showing that
      ``fraction of anatomy hidden'' does not order the results.
\item A subgroup audit, and a reproduction of the headline result to $10^{-5}$
      from public weights on different hardware and precision.
\end{itemize}

\section{Setup}

\paragraph{Data.} FairVision glaucoma OCT volumes. Downstream evaluation uses a
fixed test split of $N{=}3000$ volumes (1466 positive / 1534 negative), seed 42,
loader order fixed. Labels are byte-identical across arms, so any two arms are
paired-comparable on the same cases.

\paragraph{Pretraining.} Patch-level I-JEPA, ViT-Base, $256{\times}256$ crops,
patch size 16 (a $16{\times}16{=}256$ token grid), predictor depth 6.
Every arm resumes from the same epoch-25 checkpoint and continues to epoch 100
with identical optimiser, schedule, and effective batch size. Only the mask
sampler differs.

\paragraph{Evaluation.} The encoder is frozen. Slice features are mean-pooled,
followed by LayerNorm and a linear head; the head is selected on validation AUC
and reported on the held-out test split. This protocol is identical for all
arms and all epochs.

\subsection{Masking policies}

\begin{description}
\item[random] Stock I-JEPA multiblock: four uniform-random rectangles. No image
  content is consulted. This is the null.
\item[oracle] A band whose vertical position is located per slice by a
  per-column intensity-weighted row centroid, smoothed across columns. It uses
  a first-order intensity statistic only and consults \textbf{no segmentation
  model}.
\item[envelope] The same rectangles as \textsc{random}, but \emph{placed} on a
  retinal envelope predicted by a trained segmentation model (MIRAGE). Shape,
  size, and count are unchanged; only location changes.
\item[anatomy-v1 / anatomy-v2] Ragged connected components grown on the
  segmenter's score map, so target \emph{shape} follows tissue. v2 adds
  diagonal bridging.
\item[cover] Targets chosen to cover the anatomy while leaving a fixed visible
  fraction ($f{=}0.21$) of tissue in the context.
\end{description}

\section{Results}

\subsection{Anatomy-guided masking does not win}

Table~\ref{tab:main} gives frozen-probe test AUC at matched epochs.

\begin{table}[t]
\centering
\caption{Frozen MeanPool probe, FairVision test ($N{=}3000$). All arms share one
epoch-25 ancestor and one evaluation protocol. $\Delta$ is versus \textsc{random}
at the same epoch.}
\label{tab:main}
\small
\begin{tabular}{llccc}
\toprule
arm & guidance signal & ep50 & ep75 & ep100 \\
\midrule
random      & none (null)                & 0.8641 & 0.8723 & 0.8746 \\
oracle      & intensity centroid, no segmenter & 0.8740 & 0.8836 & \textbf{0.8855} \\
envelope    & trained segmenter          & 0.8761 & 0.8803 & 0.8807 \\
anatomy-v2  & trained segmenter (shape)  & 0.8654 & --     & --     \\
cover $f{=}0.21$ & trained segmenter (coverage) & 0.8643 & -- & -- \\
\midrule
\multicolumn{2}{l}{$\Delta$ oracle $-$ random} & $+0.0099$ & $+0.0113$ & $+0.0109$ \\
\multicolumn{2}{l}{$\Delta$ oracle $-$ envelope} & $-0.0020$ & $+0.0033$ & $+0.0047$ \\
\bottomrule
\end{tabular}
\end{table}

The oracle band gives $+0.0109$ over the null at epoch 100 (paired bootstrap
over test volumes, 95\% CI $[+0.0058,+0.0162]$, $p<0.0005$), and its margin is
stable across epochs ($+0.0099$, $+0.0113$, $+0.0109$). The segmenter-guided
\textsc{envelope} arm also beats the null, but by a margin that \emph{shrinks}
with training ($+0.0120 \to +0.0080 \to +0.0062$), and it is overtaken by the
segmentation-free oracle by epoch 75.

\subsection{Anatomical precision is not the operative variable}

We measured mask geometry directly, running each production sampler on 600
held-out slices (Table~\ref{tab:geom}).

\begin{table}[t]
\centering
\caption{Measured mask geometry, 600 slices, production samplers. ``purity'' is
the fraction of masked cells lying on tissue.}
\label{tab:geom}
\small
\begin{tabular}{lccccc}
\toprule
arm & anatomy hidden & purity & mask ratio & context kept & AUC \\
\midrule
random          & 52.2\% & 31.6\% & 43.7\% & 42.1\% & 0.8746 \\
oracle          & 62.2\% & 41.1\% & 40.0\% & 45.6\% & \textbf{0.8855} \\
envelope        & 76.9\% & 43.5\% & 46.4\% & 40.5\% & 0.8807 \\
cover $f{=}0.21$ & 74.1\% & 45.3\% & 43.3\% & 43.5\% & 0.8643 \\
anatomy-v2      & 80.3\% & \textbf{97.3\%} & 21.4\% & 67.9\% & 0.8654 \\
\bottomrule
\end{tabular}
\end{table}

The ordering is not explained by anatomical targeting. The anatomy arms place
$97$--$98\%$ of masked cells on tissue and rank last; the oracle places $41\%$
on tissue and ranks first. Across the four rectangle arms, Spearman correlation
between fraction-of-anatomy-hidden and AUC is $-0.40$ --- the wrong sign --- and
no measured geometric variable is a significant predictor at this sample size.

Two mechanisms are consistent with the data and are not separated by these
experiments. First, the oracle is the most \emph{spatially consistent} policy we
measured (highest concentration of its hit map, Gini $0.400$ versus $0.316$--$0.338$
for the other rectangle arms): it repeatedly forces the encoder to represent the
same region. Second, the anatomy arms differ from the rectangle arms on several
axes at once --- they leave $68\%$ of tokens as context versus $42\%$, and they
run at roughly half the effective masking ratio --- so their deficit cannot be
attributed to target \emph{shape} alone. We therefore do not claim that
irregular target shape is harmful; that comparison is confounded by
construction and we report it as such.

\subsection{Subgroup analysis}

Joining saved test predictions to demographics in deterministic loader order
(validated by exact label reconstruction), Black participants have the lowest
AUC in every valid probe. The disparity is a property of the evaluation, not of
a single arm: it persists across masking policies and across epochs, which
suggests it is inherited from the data and the downstream task rather than
introduced by the pretraining objective.

\subsection{Reproducibility}

The epoch-100 oracle encoder and its trained probe head are public. Re-encoding
the test split on different hardware, at fp32 rather than fp16, and with a
different encoder chunk size reproduces the reported AUC to
$9.8\times10^{-6}$ (0.8854754 versus 0.8854852), a difference of 22 discordant
pairs out of 2{,}248{,}844. Probe-time numerical precision is therefore not a
material factor in any comparison reported here.

\section{Limitations}

\textbf{One pretraining run per policy.} The paired bootstrap quantifies
sampling error on a fixed test set; it does not quantify seed-to-seed variance
of pretraining. With $n{=}1$ continuation per arm, the \emph{ranking} of arms is
not established even where individual pairwise differences are significant on
this test set. This is the principal limitation of the study.

\textbf{One visible-tissue floor.} Only $f{=}0.21$ was pretrained for the
\textsc{cover} policy, so we report no dose-response over the floor and make no
claim that any particular floor is optimal.

\textbf{Confounded shape comparison.} As noted above, the anatomy arms differ
from the rectangle arms in masking ratio, context size, and the number of
predictor loss terms simultaneously. We report their results but do not use
them to support a causal claim about target shape.

\textbf{Single dataset and task.} All results are glaucoma classification on
FairVision OCT. Generality to other modalities or endpoints is untested.

\section{Conclusion}

Guiding masked-prediction targets with a trained segmentation model did not
improve downstream glaucoma classification relative to a far cheaper
segmentation-free heuristic, and increasing the anatomical precision of the mask
was associated with worse, not better, representations. For practitioners the
practical implication is concrete: an entire segmentation stage can be removed
from this pipeline without measured loss. For the field, the result is a caution
that ``mask the clinically relevant region'' is an intuition that should be
tested rather than assumed.

\bibliographystyle{plainnat}
\bibliography{references}

\end{document}
